\section{Choosing the Workflow Type}
\label{sec:workflow-type}

The \emph{Workflow Type} section is the first choice you make after
loading images. It determines:
\begin{itemize}
  \item How each frame is compared against a reference,
  \item Which solver (Local DIC or AL-DIC) runs,
  \item How often the reference frame refreshes (incremental only),
  \item \textbf{and therefore, which frames need a Region of Interest.}
\end{itemize}

\begin{tipbox}
Get this right before drawing Regions of Interest. The Region of
Interest section below shows a live hint that updates whenever you
change the workflow type, telling you exactly which 1-based frame
numbers require a region.
\end{tipbox}

\subsection{Tracking Mode}

\begin{description}[leftmargin=1.2cm,style=nextline]
  \item[Accumulative]
    Every frame is compared directly against frame 1. Best for small,
    monotonic deformation (strain $<$ 5\%, rotation $<$ 2\textdegree).
    Fast, but fails for large cumulative motion because the reference
    and current subsets become visually too different.
  \item[Incremental]
    Each frame is compared against a rolling reference frame.
    Required for rotations $>$ 5\textdegree{} or large cumulative
    strain. Displacements from consecutive segments are composed into
    cumulative output fields. Slightly slower than accumulative but
    handles arbitrary total deformation.
\end{description}

\subsection{Solver}

\begin{description}[leftmargin=1.2cm,style=nextline]
  \item[Local DIC]
    Each subset is solved independently with IC-GN. Fast, preserves
    sharp local features, but sensitive to noise. Best for high-
    quality images and small deformations.
  \item[AL-DIC]
    Augmented Lagrangian solver that couples local IC-GN with a
    global FEM regularizer. Smoother results, better strain accuracy,
    more robust in noisy or high-gradient regions. $\sim 3\times$
    slower than Local DIC. Default.
\end{description}

When AL-DIC is selected, an indented \emph{ADMM Iterations} spinbox
appears (1--10, default 3). Higher is more regularized but slower.

\subsection{Reference Update (incremental only)}

Visible only when \emph{Tracking Mode} $=$ Incremental. Controls
when the rolling reference frame refreshes:

\begin{description}[leftmargin=1.2cm,style=nextline]
  \item[Every Frame]
    Reference resets every frame. Most robust for large deformation
    (smallest per-step ratio). A Region of Interest is only needed on
    frame 1 --- pyALDIC warps it forward automatically.
  \item[Every N Frames]
    Reference resets every $N$ frames. You need a Region of Interest
    on frame 1, $N+1$, $2N+1$, \ldots. The live hint in the Region of
    Interest section lists the required frames.
  \item[Custom Frames]
    You type a comma-separated list of reference frames, e.g.\
    \texttt{1, 6, 12, 24}. Frame 1 is always included.
\end{description}

\subsection{Which combination to pick?}

\begin{center}
\begin{tabular}{lll}
  \toprule
  Experiment & Tracking & Reference Update \\
  \midrule
  Soft material, $<$ 2\% strain & Accumulative & --- \\
  Soft material, 2--30\% strain & Incremental & Every Frame \\
  Rigid rotation, $>$ 5\textdegree{} & Incremental & Every Frame \\
  Dynamic event / impact     & Incremental & Every Frame \\
  Quasi-static with known    & Incremental & Every N or Custom \\
  \quad load steps           &            & (align $N$ with steps) \\
  \bottomrule
\end{tabular}
\end{center}

\begin{notebox}
When in doubt, use \emph{Incremental + Every Frame}. It is slightly
slower than accumulative but handles every deformation magnitude
well, and only needs one Region of Interest drawn on frame 1.
\end{notebox}
